\documentclass[11pt]{article}
\usepackage[margin=1in]{geometry}
\usepackage{graphicx}
\usepackage{booktabs}
\usepackage{amsmath,amssymb}
\usepackage{caption}
\usepackage{hyperref}
\usepackage{natbib}
\usepackage{times}

\title{Glutamate neurotransmission from leptin-receptor neurons of the ventral premammillary nucleus is required for female sexual maturation and reproduction}
\author{}
\date{}

\begin{document}
\maketitle

\begin{abstract}
The adipocyte hormone leptin acts as a permissive metabolic signal for the hypothalamic--pituitary--gonadal axis, and leptin-receptor (LepRb) neurons of the ventral premammillary nucleus (PMv) have been identified as a privileged relay between energy stores and reproduction. However, the neurotransmitter through which these neurons convey the leptin signal has not been established. Because most PMv LepRb neurons are glutamatergic (Vglut2$^{+}$), we hypothesised that glutamate is the indispensable effector. We combined chemogenetic activation of PMv LepRb neurons with automated serial blood sampling for luteinising hormone (LH), conditional deletion of Vglut2 from LepRb cells (LepRb$\Delta$Vglut2), and simultaneous viral re-expression of Lepr together with deletion of Vglut2 selectively in the PMv of LepR-null (LepR$^{loxTB}$) mice. Chemogenetic activation triggered LH release within 10--20 min in 55\% (8/15) of CNO-injected PMv-hit animals and 100\% (4/4) of clozapine-injected PMv-hit animals; LH peaks correlated tightly with PMv cFos counts (Pearson $r=0.67$, $p=0.0007$). LepRb$\Delta$Vglut2 females developed late-onset obesity and showed prolonged estrous cycles ($3.31\pm0.36$ vs $4.47\pm0.30$ cycles per 30 d; $t_{26}=2.45$, $p=0.021$), reduced arcuate GnRH immunoreactivity ($t_{5}=8.01$, $p=0.0005$), blunted LH responses to kisspeptin-10 (AUC $t_{11}=3.14$, $p=0.009$), and $\sim$50\% lower hypothalamic Tac2 expression ($t_{8}=3.39$, $p=0.009$). Viral re-expression of Lepr in the PMv restored leptin-induced pSTAT3 in both LepR$^{loxTB}$ and LepR$^{loxTB}$;Vglut2$^{flox}$ mice, but rescue of puberty and fertility occurred only when glutamate transmission was preserved: correctly PMv-injected LepR$^{loxTB}$;Vglut2$^{flox}$ females showed no vaginal opening and did not become pregnant. Glutamate is therefore the indispensable neurotransmitter through which PMv LepRb neurons relay the permissive signal of leptin to the reproductive axis.
\end{abstract}

\section{Introduction}

The capacity of a female mammal to enter puberty and sustain adult reproduction depends on whether body energy stores are adequate to underwrite the metabolic cost of gestation and lactation. The adipocyte-derived hormone leptin communicates this information to the brain and acts as a permissive signal for the reproductive axis: mice and humans lacking leptin or its receptor are obese and fail to enter puberty, and exogenous leptin reverses hypogonadotropic hypogonadism in leptin-deficient subjects~\citep{ahima1997leptin,farooqi2014twenty}. Leptin acts through its long-form receptor, LepRb, on neurons distributed across many hypothalamic and brainstem nuclei; however, the cellular identity and neurochemical effectors through which the signal reaches gonadotropin-releasing hormone (GnRH) neurons have remained incompletely resolved~\citep{louis2011molecular}.

GnRH neurons do not themselves express LepRb, so leptin must act through intermediary populations that gate GnRH output. Candidate relays include Kiss1/neurokinin B/dynorphin (KNDy) neurons of the arcuate nucleus (Arc), proopiomelanocortin (POMC) neurons of the Arc, and LepRb neurons of the ventral premammillary nucleus (PMv). Yet genetic deletion of LepRb specifically from Kiss1 neurons does not disrupt puberty or fertility, and only a small minority of Arc KNDy neurons coexpress LepRb, arguing against a direct kisspeptin-centred model of leptin action~\citep{donato2011leptin,cravo2011kiss1,louis2011molecular,true2011leptin}. A melanocortin relay in which Arc POMC neurons signal onto KNDy cells has been proposed to convey leptin's permissive effect on puberty, consistent with direct LepRb action on POMC cells in other contexts~\citep{manfredilozano2016defining,berglund2012direct}. In contrast, earlier work showed that bilateral lesions of the PMv blunt leptin-induced sexual maturation and that unilateral endogenous re-expression of LepRb in the PMv alone can restore puberty and fertility in LepR-null females, pinpointing the PMv as a privileged site of leptin action~\citep{donato2011leptin,mahany2018obesity}.

PMv LepRb neurons have the signatures expected of a metabolic-to-reproductive relay. Electrophysiology shows that leptin depolarises a majority of them through a PI3K/TRPC-dependent mechanism~\citep{williams2011acute}. Anatomically, they monosynaptically innervate GnRH neurons and are activated by both leptin and sexual odorants~\citep{leshan2009direct}. Notably, most PMv LepRb neurons are glutamatergic (Vglut2$^{+}$), raising the largely untested hypothesis that glutamate is the neurotransmitter through which leptin's permissive signal is conveyed. Circumstantial support comes from transmitter-specific LepR deletions: selective removal of LepRb from GABAergic---but not glutamatergic---populations produces obesity with delayed puberty, showing that the transmitter identity of LepRb neurons matters for reproductive physiology and may segregate across sites~\citep{martin2014leptinresponsive}.

Despite this convergent circumstantial evidence, no study has directly tested (i) whether acute activation of PMv LepRb neurons is sufficient to trigger GnRH/LH release, (ii) whether conditional deletion of Vglut2 from LepRb cells impairs puberty and adult reproduction, or (iii) whether the requirement for glutamate is specifically local to the PMv. We address these gaps with three complementary strategies. First, we use chemogenetic (DREADD hM3Dq) stimulation to test whether activation of PMv LepRb neurons suffices to drive LH release on a physiologically relevant timescale. Second, we delete Vglut2 from all LepRb cells to ask whether glutamate released from LepRb neurons is required for pubertal maturation, cyclicity, pituitary responsiveness to kisspeptin, and fertility. Third, we re-express Lepr in the PMv of LepR-null (LepR$^{loxTB}$) females while simultaneously deleting Vglut2 in the same cells, to test whether glutamate is specifically the permissive transmitter at the PMv node. Together these experiments show that glutamate is an indispensable effector of the PMv leptin-to-reproduction pathway.

\section{Methods}

\subsection{Animals}
All experiments used adult female mice of the following genotypes, maintained on a congenic C57BL/6J background: LepRb-Cre (for chemogenetic activation of LepRb neurons, $n=19$ CNO and $n=4$ clozapine recipients with PMv-targeted AAV); LepRb-Cre;Vglut2$^{flox}$ crossed to generate LepRb$\Delta$Vglut2 experimentals and LepRb-Cre or Vglut2$^{flox}$ littermate controls; LepRb-Cre;Kiss1-hrGFP reporter mice ($n=4$ PMv-hits) for visualising Kiss1 neurons; LepR$^{loxTB}$ LepR-null mice carrying a loxP-flanked transcription-blocking cassette between exons 16 and 17 of \textit{Lepr}~\citep{berglund2012direct}; and LepR$^{loxTB}$;Vglut2$^{flox}$ compound mutants. Wild-type mice ($n=13$; six with AAV-GFP, seven uninjected) served as fertile controls and uninjected LepR$^{loxTB}$ mice ($n=4$) as infertile controls for the rescue experiment.

\subsection{Stereotaxic surgeries}
Adult females (8--10 weeks old) were anaesthetised with isoflurane and placed in a stereotaxic frame. PMv coordinates, measured from the rostral rhinal vein, were anteroposterior $-5.4$~mm; mediolateral $\pm0.5$~mm from the exposed superior sagittal sinus; and dorsoventral $-5.4$~mm from the dura mater. $\sim$50~nL of AAV-hM3Dq-mCherry (DREADD activation cohort) or AAV-Cre (rescue cohort) was injected unilaterally. Mice were monitored twice daily for 3 days and daily for 7 further days. Experiments were performed 1 month after surgery.

\subsection{Chemogenetic activation and LH sampling}
Three days after surgery, LepRb-Cre females ($n=19$ for CNO, $n=4$ for clozapine) were connected to the Culex automated blood-sampling system. The following day, three baseline samples were drawn every 10~min to establish pre-injection LH levels, followed by an intravenous injection of either clozapine-$N$-oxide (CNO) or clozapine at 0.5~mg/kg in 50~$\mu$L at time 0. Blood samples were then drawn every 10~min for 1~h. An additional control cohort of LepRb-Cre females with no AAV injection ($n=7$) received iv clozapine alone to control for non-specific ligand effects. Two mice with AAVs landing outside the PMv were classified as ``PMv-misses'' ($n=2$), and a further group of AAV-mCherry-injected animals ($n=5$) served as viral controls. Two hours after injection animals were perfused for cFos and mCherry immunolabelling in the PMv and Arc.

\subsection{Kisspeptin-10 challenge}
LepRb-Cre ($n=7$) and LepRb$\Delta$Vglut2 ($n=6$) females were cannulated as above and, after baseline sampling, received an intraperitoneal injection of kisspeptin-10 at 65~$\mu$g/kg (Kisspeptin-10, \#048-56 Phoenix Pharmaceuticals)~\citep{gottsch2004role}. Blood was sampled every 7~min for 30~min. The area under the LH curve (AUC) was computed after subtracting the $-10$~min value from every timepoint as baseline 0.

\subsection{AAV-Cre re-expression of Lepr and concomitant Vglut2 deletion}
To ask whether glutamate is specifically required at the PMv, we used the LepR$^{loxTB}$ model, in which LepR signalling can be restored by Cre-dependent excision of the transcription-blocking cassette. LepR$^{loxTB}$ females received unilateral PMv injections of AAV-Cre as a positive control ($n=5$); LepR$^{loxTB}$;Vglut2$^{flox}$ females ($n=9$) received the same injection, which simultaneously re-expressed Lepr and deleted Vglut2 only in Cre-targeted PMv neurons. Six of these females had AAV-Cre centred in the PMv; two had injections landing in adjacent dorsomedial or ventromedial hypothalamic nuclei and were excluded. Six to eight weeks after surgery, mice were injected ip with leptin (3~mg/kg) 1.5~h before perfusion for pSTAT3 immunohistochemistry.

\subsection{Histology, in situ hybridisation and qPCR}
Mice were perfused with 0.1~M phosphate-buffered saline (PBS) followed by 10\% neutral-buffered formalin; brains were cryoprotected in 20\% sucrose and cut on a freezing microtome into four series of 30-$\mu$m coronal sections. Reproductive organs were paraffin-embedded and stained with haematoxylin and eosin. Immunofluorescence used 3\% normal donkey serum with 0.25\% Triton X-100 blocking and standard primary-antibody protocols. For pSTAT3, endogenous peroxidases were blocked with 30\% H$_{2}$O$_{2}$, followed by 0.3\% glycine and 0.03\% SDS; detection used biotinylated secondary antibody and diaminobenzidine (ABC Kit, Vector Laboratories). Radioactive in situ hybridisation for Vglut2 used a T7-transcribed $^{35}$S-labelled antisense riboprobe hybridised overnight at $57^{\circ}$C in 50\% formamide. Chromogenic duplex in situ hybridisation for Lepr (Green) and Vglut2 (Fast Red) used BaseScope Duplex reagents (ACD, \#323810) on 16-$\mu$m fresh-frozen sections, with 50\% Gill's haematoxylin counterstain. qPCR was performed on 20~ng cDNA in triplicate 10-$\mu$L reactions using SYBR Green Taq MasterMix on a CFX-384 instrument (Bio-Rad); pituitary genes were normalised to \textit{Actb} and mediobasal hypothalamic genes to \textit{B2m}.

\subsection{LH assay}
LH was measured by an in-house sandwich ELISA. Functional sensitivity was 0.016~ng/mL and the reportable range was 0.016--4.0~ng/mL; intra-assay coefficient of variation was 6.24\% and the intra-assay CV was 2.2\% for repeated measures; inter-assay CVs were 7.3\% (low QC, 0.13~ng/mL), 5.0\% (medium QC, 0.8~ng/mL) and 6.5\% (high QC, 2.3~ng/mL).

\subsection{Image analysis and statistics}
Photomicrographs were acquired on an Axio Imager M2 (Carl Zeiss Microscopy). To quantify GnRH immunoreactivity in the Arc, a 250~$\mu$m wide $\times$ 140~$\mu$m high rectangle was placed on the 3V wall, avoiding the saturated signal of the tuberoinfundibular sulcus, and integrated density (mean pixel intensity $\times$ area) was calculated. Data are expressed as mean $\pm$ SEM. One-sample $t$-tests were used for LH AUC against a baseline of 0; two-group comparisons used unpaired Student's $t$-tests; multi-group comparisons used one-way ANOVA with post hoc tests; non-normal distributions were analysed with the Kruskal--Wallis test; repeated-measures LH data were analysed with mixed-effects models followed by Sidak post hoc tests; associations used Pearson correlation.

\section{Results}

\subsection{Chemogenetic activation of PMv LepRb neurons induces LH release}
To test whether acute activation of PMv LepRb neurons is sufficient to drive LH release in fed adult females, we unilaterally injected AAV-hM3Dq-mCherry into the PMv of LepRb-Cre mice and, four weeks later, connected the animals to the Culex automated blood-sampling system for serial LH measurement (Fig.~\ref{fig:design}). Animals were classified post hoc as ``PMv-hits'' when mCherry expression was centred in the PMv, ``PMv-misses'' when the injection landed outside the PMv ($n=2$), or viral controls (AAV-mCherry, $n=5$). A separate cohort of LepRb-Cre females with no AAV injection ($n=7$) received intravenous clozapine alone to control for clozapine-independent effects.

Intravenous CNO (0.5~mg/kg) induced a rise of serum LH above 0.5~ng/mL within 10--20~min in 8 of 15 PMv-hit animals ($\sim$55\%), with peak LH levels between 0.65 and 3.11~ng/mL. One-sample $t$-tests against a baseline of zero confirmed the response: PMv-hit animals with an LH rise yielded $t_{7}=3.54$, $p=0.009$, whereas PMv-hit animals without an LH rise ($t_{6}=1.37$, $p=0.22$) and PMv-miss plus AAV-mCherry controls ($t_{6}=1.00$, $p=0.36$) did not. In PMv-miss animals, one of two mice showed a modest LH rise (0.57~ng/mL) 10~min after CNO and the other remained below the limit of detection. Of five AAV-mCherry-injected controls only one showed a rise (0.93~ng/mL), and at 30~min rather than the 10--20-min window (Fig.~\ref{fig:responders}). Because CNO can be back-metabolised to clozapine, we also tested clozapine directly in a cohort expressing AAV-hM3Dq in the PMv ($n=4$): all four animals showed an LH rise within 10~min (0.87--1.18~ng/mL; $t_{3}=8.79$, $p=0.003$), whereas only 2 of 7 (28.5\%) no-AAV animals given clozapine showed any rise ($t_{6}=1.83$, $p=0.12$).

Activation mapping with cFos immunofluorescence was consistent with these outcomes. The number of cFos-immunoreactive neurons per section in the PMv differed markedly across groups (one-way ANOVA, $F_{2,19}=81.66$, $p<0.0001$), with PMv-hit animals showing the highest counts and negative controls the lowest. For the clozapine cohort, PMv cFos was again higher in hM3Dq hits than in no-AAV controls ($t_{8}=17.46$, $p<0.0001$). Across all animals, the peak LH level correlated strongly with the number of PMv cFos neurons (Pearson $r=0.67$, $p=0.0007$), indicating that the magnitude of PMv LepRb recruitment quantitatively predicts the LH response. Together these data show that acute activation of PMv LepRb neurons is sufficient to drive LH secretion and that only a sufficiently large or specific subset of PMv LepRb neurons must be recruited to cross the threshold for LH release: the lower response rate with CNO (55\%) compared with clozapine (100\%) is consistent with a ligand-potency or dose-threshold effect at the hM3Dq receptor.

\subsection{Viral contamination of the posterior arcuate produces inconsistent LH secretion}
Because stereotaxic AAV injections in the PMv can spread into adjacent nuclei, we next asked whether spillover into the posterior arcuate nucleus---home to POMC and KNDy neurons---might contribute to the LH response. Peak LH correlated with the number of cFos plus mCherry double-labelled cells in the posterior Arc (Pearson $r=0.62$, $p=0.002$) but not with the total number of cFos cells in the Arc alone ($r=0.3$, $p=0.16$), indicating that activation of Arc LepRb-Cre neurons (mCherry$^{+}$) directly activated by the virus was better correlated with LH than bulk Arc activity. A Kruskal--Wallis comparison of Arc cFos$^{+}$mCherry$^{+}$ cells across groups confirmed a group effect ($H=14.63$, $p<0.0001$).

We then quantified how many of the activated Arc neurons were POMC or Kiss1. In PMv-hit animals, POMC immunoreactivity was observed in $9.3\pm2.6\%$ of Arc cFos-positive neurons (corresponding to $12.3\pm5.0$ cells per section), and in $9.6\pm3.1\%$ of cFos neurons in hM3Dq clozapine-treated animals ($2.3\pm1.3$ cells). Of the POMC/cFos double-labelled neurons, about half also coexpressed mCherry in the PMv-hit group ($7.3\pm3.2$ cells) but essentially none in the hM3Dq clozapine group ($0.3\pm0.3$ cells). In LepRb-Cre;Kiss1-hrGFP PMv-hit animals ($n=4$), only $2.5\pm0.95\%$ of Arc cFos cells were Kiss1-hrGFP$^{+}$, corresponding to $1.8\pm0.3$ cells per section; most of these also expressed mCherry, suggesting that viral contamination reached at most one or two KNDy neurons per animal.

These observations argue that spillover activation of POMC or Kiss1 neurons contributed minimally, if at all, to the LH response: the absolute numbers of virally activated POMC/Kiss1 neurons are very small, and cFos/LH increases were driven overwhelmingly by PMv activation. The partial correlation with double-labelled Arc cells most likely reflects the fact that the same animals with large PMv hits also had larger spread into adjacent territories, rather than implying that Arc POMC or KNDy activation is the primary effector. The consistently higher response rate in AAV-hM3Dq PMv-hits given CNO (55\%) or clozapine (100\%) versus PMv-misses, AAV-mCherry and no-AAV controls (0\% and 28.5\%), together with strong PMv cFos only in responders, makes it unlikely that the observed LH rises represent spontaneous pulsatile secretion.

\subsection{LepRb$\Delta$Vglut2 females develop late-onset obesity and impaired reproduction}
Having established that PMv LepRb activation can drive LH, we next asked whether glutamate released from LepRb cells is required for reproductive function. We crossed LepRb-Cre with Vglut2$^{flox}$ mice to delete the vesicular glutamate transporter 2 conditionally from LepRb-expressing cells. Adult LepRb$\Delta$Vglut2 mice (13--26 weeks) developed a late-onset obesity phenotype (Fig.~\ref{fig:cycles}): males showed a significant main effect of genotype on body mass (two-way ANOVA, $F_{1,10}=13.1$, $p=0.0047$; post hoc comparisons: week 20, $t_{140}=3.27$, $p=0.019$; week 26, $t_{140}=5.47$, $p<0.0001$), and females showed a trend-level main effect ($F_{1,11}=4.34$, $p=0.06$) with significant post hoc differences at weeks 23, 25 and 26 ($t_{154}=4.05$, $p=0.0011$; $t_{154}=3.1$, $p=0.03$; $t_{154}=5.59$, $p<0.0001$). The weight gain reflected increased fat mass in both sexes (males: $t_{9}=3.77$, $p=0.004$; females: $t_{11}=3.7$, $p=0.0035$), reproducing the metabolic phenotype previously reported for this cross~\citep{martin2014leptinresponsive}.

We then assessed whether the loss of Vglut2 in LepRb cells altered pubertal maturation. The age of vaginal opening did not differ between controls and LepRb$\Delta$Vglut2 females ($t_{18}=0.84$, $p=0.41$) and body mass at vaginal opening was similarly unaffected ($t_{18}=1.47$, $p=0.16$). In contrast, the age of first estrus was delayed in LepRb$\Delta$Vglut2 mice ($t_{14}=2.45$, $p=0.028$) with no corresponding difference in body mass ($t_{14}=0.59$, $p=0.56$). In males, the timing of complete balanopreputial separation was unchanged ($t_{11}=1.80$, $p=0.1$) although body mass at that milestone was slightly higher in mutants ($t_{11}=2.43$, $p=0.033$). Thus, the pubertal phenotype appeared predominantly in the estrous cycle rather than in the first morphological markers of puberty.

Adult cyclicity was substantially impaired. Cycle duration was longer in LepRb$\Delta$Vglut2 females ($t_{25}=3.71$, $p=0.001$), with significantly more time spent in proestrus ($t_{26}=2.81$, $p=0.009$) and diestrus ($t_{26}=2.39$, $p=0.024$) and a trend toward less time in estrus ($t_{26}=1.97$, $p=0.059$). The number of cycles completed in 30 days was reduced from $4.47\pm0.30$ in LepRb-Cre controls to $3.31\pm0.36$ in LepRb$\Delta$Vglut2 mutants ($t_{26}=2.45$, $p=0.021$; Fig.~\ref{fig:cycles}). Fertility was also compromised: 75\% of LepRb$\Delta$Vglut2 females produced at least one litter within 60 days of pairing with proven male breeders, compared with 100\% of controls. Among fertile mutants, latency to first litter was not different (Mann--Whitney $p=0.32$) nor was litter size ($7.50\pm1.41$ pups versus $7.86\pm1.14$; $t_{11}=0.20$, $p=0.85$). These data collectively show that loss of Vglut2 from LepRb cells produces a reproductive phenotype that appears before and beyond the metabolic phenotype: cycle length and first estrus are altered in LepRb$\Delta$Vglut2 females between 10 and 16 weeks of age, well before the late-onset divergence in body weight. Having shown that glutamate from LepRb cells is required for normal cyclicity, we next asked at which level of the reproductive axis the lesion occurs.

\subsection{Reduced GnRH-ir and blunted pituitary response to kisspeptin in LepRb$\Delta$Vglut2 females}
To locate the lesion in the reproductive axis, we examined hypothalamic GnRH content and pituitary responsiveness. GnRH immunoreactivity in the Arc, quantified as integrated density in a 250~$\mu$m $\times$ 140~$\mu$m region on the 3V wall, was markedly reduced in LepRb$\Delta$Vglut2 females ($n=3$) compared with Vglut2$^{flox}$ controls ($n=4$; $t_{5}=8.01$, $p=0.0005$). Basal LH in diestrus showed a trend toward reduction ($t_{10}=1.91$, $p=0.08$) but did not reach significance, consistent with a deficit in stimulated rather than tonic release.

We next challenged animals with exogenous kisspeptin-10 (65~$\mu$g/kg ip) to probe pituitary responsiveness downstream of KNDy neurons~\citep{gottsch2004role}. LH levels were higher in control than in LepRb$\Delta$Vglut2 females after the injection (mixed-effects model: time $F_{1.14,12.19}=15.91$, $p=0.001$; genotype $F_{1,11}=8.26$, $p=0.015$; time $\times$ genotype $F_{3,32}=5.00$, $p=0.006$). Sidak post hoc tests showed genotype differences at 7~min ($t_{10.97}=3.79$, $p=0.01$) and 14~min ($t_{9.4}=3.22$, $p=0.04$). The AUC of LH over the 30-min window was significantly smaller in LepRb$\Delta$Vglut2 animals ($t_{11}=3.14$, $p=0.009$), identifying a deficit in the acute pituitary response to kisspeptin.

To determine whether this deficit reflected changes in gonadotrope machinery or in the KNDy network upstream, we performed qPCR on pituitary and mediobasal hypothalamic (MBH) tissue from diestrus females. Pituitary expression of \textit{Gnrhr} ($t_{10}=0.86$, $p=0.41$), \textit{Lh$\beta$} ($t_{10}=1.47$, $p=0.17$) and \textit{Fsh$\beta$} ($t_{9}=1.63$, $p=0.14$) did not differ between genotypes, and the common $\alpha$-subunit \textit{Cga} was likewise unchanged (Vglut2$^{flox}$: $100\pm17.47\%$; LepRb$\Delta$Vglut2: $245.3\pm89.91\%$; $t_{8}=1.20$, $p=0.26$). In the MBH, \textit{Kiss1} ($t_{8}=1.42$, $p=0.19$), \textit{Pdyn} ($t_{8}=0.02$, $p=0.98$), \textit{Kiss1r} ($t_{8}=1.44$, $p=0.19$) and \textit{Tac3r} ($t_{8}=0.99$, $p=0.35$) were unchanged, but \textit{Tac2} was reduced by approximately 50\% in LepRb$\Delta$Vglut2 females ($t_{8}=3.39$, $p=0.009$; Fig.~\ref{fig:mbhqpcr}). The dissociation between unchanged pituitary gene expression and a blunted LH response to exogenous kisspeptin-10, coupled with reduced arcuate GnRH content, points to a lesion at the level of GnRH terminals or of the KNDy autofeedback loop that shapes pulsatile GnRH release~\citep{navarro2011regulation,navarro2013interactions}.

\subsection{LepR re-expression in the PMv requires glutamate for pubertal rescue and fertility}
To test whether glutamate is specifically required at the PMv, we exploited the LepR$^{loxTB}$ null allele, in which Cre excises a transcription-blocking cassette and restores Lepr expression~\citep{berglund2012direct}. We injected AAV-Cre unilaterally into the PMv of LepR$^{loxTB}$ ($n=5$, positive control) and LepR$^{loxTB}$;Vglut2$^{flox}$ ($n=9$) females; in the compound-mutant line, the same injection simultaneously restored LepR signalling and deleted Vglut2 in Cre-targeted neurons. Six of the nine compound-mutant females received AAV-Cre centred in the PMv and were analysed; two with injections landing in adjacent dorsomedial or ventromedial hypothalamic nuclei were excluded from analysis. Wild-type mice ($n=13$; six with AAV-GFP, seven uninjected) and uninjected LepR$^{loxTB}$ mice ($n=4$) served as fertile and infertile controls respectively.

Quantification of leptin-induced pSTAT3 immunoreactivity revealed that viral re-expression of Lepr produced similar numbers of PMv pSTAT3-positive neurons in AAV-Cre-injected LepR$^{loxTB}$ and AAV-Cre-injected LepR$^{loxTB}$;Vglut2$^{flox}$ mice, in both cases comparable to wild-type animals and far greater than uninjected LepR$^{loxTB}$ controls (one-way ANOVA, $F_{3,13}=73.53$, $p<0.0001$; Fig.~\ref{fig:rescue}). In the arcuate nucleus, where LepRb$^+$ neurons are not directly manipulated by the viral injection, overall group effects on pSTAT3 were also apparent ($F_{3,16}=171.8$, $p<0.0001$). Chromogenic in situ hybridisation confirmed that Vglut2 signal in the PMv was decreased only in the compound-mutant line when the injection was on target ($t_{8}=4.20$, $p=0.003$), showing successful simultaneous re-expression of Lepr and ablation of Vglut2.

The behavioural and reproductive outcome was unambiguous. Correctly PMv-injected LepR$^{loxTB}$;Vglut2$^{flox}$ females showed no sign of puberty onset (no vaginal opening) and no pregnancy by the end of the experiment ($\sim$60 days), in contrast to the rescue obtained in AAV-Cre-injected LepR$^{loxTB}$ mice. Quantification of corpora lutea in the ovary, an index of recent ovulation, showed a pronounced group effect ($F_{3,17}=31.16$, $p<0.0001$), with AAV-Cre LepR$^{loxTB}$ rescues producing corpora lutea at wild-type frequencies but the compound-mutant rescues failing to ovulate. Because restoration of PMv LepR signalling alone was insufficient when glutamate was simultaneously removed from those neurons, these data establish that glutamate is the indispensable neurotransmitter through which PMv LepRb neurons relay the leptin signal to the downstream reproductive network (Fig.~\ref{fig:model}).

\section{Discussion}

Our study resolves a long-standing question about how metabolic status gates reproduction: glutamate released by LepRb neurons of the ventral premammillary nucleus is indispensable for pubertal maturation, adult cyclicity and fertility. Three lines of evidence converge on this conclusion. First, chemogenetic activation of PMv LepRb neurons is sufficient to trigger LH release within 10--20~min, with LH peaks tightly correlated to PMv cFos counts (Pearson $r=0.67$, $p=0.0007$). Second, conditional deletion of Vglut2 from LepRb cells produces delayed first estrus, prolonged estrous cycles, reduced Arc GnRH immunoreactivity, a blunted LH response to exogenous kisspeptin-10, and $\sim$50\% lower hypothalamic Tac2 expression. Third, simultaneous re-expression of Lepr and deletion of Vglut2 in PMv neurons of LepR-null mice restores LepR-driven pSTAT3 but fails to rescue puberty or fertility, whereas re-expression of Lepr alone does rescue. The loss- and gain-of-function experiments therefore meet at the same node: glutamate from PMv LepRb neurons is the permissive effector.

These findings substantially refine the prevailing picture of how leptin permits reproduction. Classic work established the PMv as a leptin-sensitive relay sufficient for reproductive rescue~\citep{donato2011leptin,mahany2018obesity,leshan2009direct,williams2011acute}, while complementary studies revealed that Kiss1 neurons largely lack LepRb and that removal of LepRb from Kiss1 neurons does not disrupt fertility~\citep{donato2011leptin,cravo2011kiss1,louis2011molecular,true2011leptin}. Our data are consistent with both observations and add the missing biochemical piece: leptin's permissive action in the PMv is transduced not by a peptidergic co-transmitter but by classical glutamatergic neurotransmission onto downstream partners, likely including GnRH neurons directly and the KNDy network indirectly. The reduction in Tac2 expression in LepRb$\Delta$Vglut2 mice and the blunted LH response to exogenous kisspeptin-10---despite intact pituitary gene expression---suggest that loss of glutamatergic tone from PMv LepRb neurons disrupts KNDy-dependent pulsatile GnRH output~\citep{navarro2011regulation,navarro2013interactions,gottsch2004role}.

The observation that only $\sim$55\% of PMv-hit mice responded to CNO but 100\% of PMv-hit mice responded to clozapine argues for a dose- or potency-threshold effect rather than for anatomical heterogeneity alone. Because CNO must cross the blood--brain barrier and, in part, be back-metabolised to clozapine to activate hM3Dq, effective ligand concentrations at the DREADD receptor differ between the two regimens. Consistent with a threshold, 2 of 7 (28.5\%) no-AAV animals given clozapine showed an LH rise, indicating that clozapine alone is not strictly silent in this paradigm, yet the proportion of responders was markedly lower than in hM3Dq-expressing animals. The tight correlation between PMv cFos and peak LH across animals further supports a graded, recruitment-dependent LH response. This result has two practical implications: PMv-targeted DREADD studies should include direct clozapine controls in addition to CNO, and the number or subset of PMv LepRb neurons that must be recruited to release LH is a quantitative parameter that shapes both the magnitude and incidence of the response.

Contamination of neighbouring nuclei by the AAV is an unavoidable concern with stereotaxic injections, but several features of our data argue that PMv activation, and not spillover, drives LH secretion. Arc POMC coexpression was observed in only $\sim$9\% of Arc cFos cells, corresponding to at most $\sim$7 virally targeted POMC/cFos neurons per animal, and a similar fraction appeared in the clozapine cohort where no such cells expressed mCherry. In Kiss1-hrGFP reporters, only $\sim$2.5\% of Arc cFos cells were Kiss1-positive, corresponding to 1--2 KNDy neurons per animal. Although Arc cFos/mCherry counts correlated with peak LH, cFos alone did not, suggesting that the correlation reflects co-variation between the magnitude of the PMv hit and incidental spread rather than an independent Arc mechanism. These considerations make it unlikely that direct POMC~\citep{manfredilozano2016defining,berglund2012direct} or KNDy activation~\citep{cravo2011kiss1} in the Arc accounts for the observed LH response. Having thus excluded the most obvious confounds of AAV spread, we next turned to a genetic loss-of-function strategy to test whether glutamate released from LepRb cells is itself required for reproduction.

Our work has limitations. The chemogenetic experiments used unilateral injections, which limit the number of recruited PMv LepRb neurons; this may underestimate the system's capacity for LH release. The LepRb$\Delta$Vglut2 cross removes Vglut2 from all LepRb cells, not only those in the PMv, so extra-PMv LepRb glutamatergic populations may contribute to the reproductive phenotype; however, the third experiment---simultaneous PMv-restricted LepR re-expression and Vglut2 deletion in LepR-null mice---argues that the PMv contribution is both necessary and sufficient on its own. We did not examine the male phenotype in depth beyond body-mass and balanopreputial-separation measurements, nor did we perform cell-type-resolved recordings of glutamatergic output from identified PMv LepRb neurons onto GnRH or KNDy cells. The blunted Tac2 expression and kisspeptin-10 response are suggestive of a functional KNDy lesion, but direct electrophysiological evidence for monosynaptic glutamatergic input from PMv LepRb neurons to KNDy cells remains to be established.

Several avenues flow naturally from these results. Selective deletion of Vglut2 in the PMv, using AAV-Cre in Vglut2$^{flox}$ mice without concurrent LepR manipulation, would test necessity of PMv glutamate in metabolically intact animals and disentangle the metabolic and reproductive deficits. Single-cell transcriptomics of PMv LepRb neurons would reveal whether subsets differ in cotransmitter, projection target, or leptin responsiveness, potentially explaining why only some CNO-treated hits release LH. Finally, the permissive role of PMv glutamate identifies the glutamatergic synapse between PMv LepRb neurons and GnRH or KNDy cells as a candidate node where metabolic dysfunction of obesity or eating disorders could translate into altered pubertal timing and reproductive outcomes in humans~\citep{farooqi2014twenty,ahima1997leptin,amstalden2011neuroendocrine}. Together with earlier anatomical and molecular work, our results place glutamate at the heart of the metabolic control of female reproduction and establish a tractable circuit for future mechanistic and translational study.

\section*{Data Availability}
Source data supporting the findings of this study are provided in the experimental log accompanying this manuscript. Further details are available from the corresponding author upon reasonable request.

\bibliographystyle{unsrtnat}
\bibliography{references}

% ==========================================================================
% Tables
% ==========================================================================
\begin{table}[htbp]
\centering
\caption{Chemogenetic activation of PMv LepRb neurons: LH response by group.}
\label{tab:dreaddsummary}
\begin{tabular}{lcccc}
\toprule
Group & $n$ & LH responders & Peak LH range (ng/mL) & PMv cFos statistic \\
\midrule
PMv-hit + CNO & 15 & 8 (55\%) & 0.65--3.11 & $F_{2,19}=81.66$, $p<0.0001$ \\
PMv-hit + clozapine & 4 & 4 (100\%) & 0.87--1.18 & $t_{8}=17.46$, $p<0.0001$ \\
PMv-miss + AAV-mCherry & 7 & 0 / 1 at 30 min & $<$0.57 / 0.93 & (see above) \\
No AAV + clozapine & 7 & 2 (28.5\%) & not reported & --- \\
\bottomrule
\end{tabular}
\end{table}

\begin{table}[htbp]
\centering
\caption{Reproductive and metabolic phenotype of LepRb$\Delta$Vglut2 mice. Values shown as means $\pm$ SEM where available.}
\label{tab:metabolic}
\begin{tabular}{p{4.2cm}lp{6.0cm}}
\toprule
Measure (Sex) & Direction in LepRb$\Delta$Vglut2 & Statistic \\
\midrule
Body mass, genotype effect (M) & increased from week 20 on & $F_{1,10}=13.1$, $p=0.0047$; week 20 $t_{140}=3.27$, $p=0.019$; week 26 $t_{140}=5.47$, $p<0.0001$ \\
Body mass, genotype effect (F) & increased from week 23 on & $F_{1,11}=4.34$, $p=0.06$; week 23 $t_{154}=4.05$, $p=0.0011$; week 26 $t_{154}=5.59$, $p<0.0001$ \\
Fat mass (M) & increased & $t_{9}=3.77$, $p=0.004$ \\
Fat mass (F) & increased & $t_{11}=3.7$, $p=0.0035$ \\
Age of vaginal opening (F) & unchanged & $t_{18}=0.84$, $p=0.41$ \\
Age of first estrus (F) & delayed & $t_{14}=2.45$, $p=0.028$ \\
Cycle duration (F) & longer & $t_{25}=3.71$, $p=0.001$ \\
Cycles per 30 days (F) & $4.47\pm0.30 \to 3.31\pm0.36$ & $t_{26}=2.45$, $p=0.021$ \\
Fertility within 60 d (F) & $100\% \to 75\%$ & Mann--Whitney latency $p=0.32$ \\
First-litter pup number (F) & $7.86\pm1.14 \to 7.50\pm1.41$ & $t_{11}=0.20$, $p=0.85$ \\
\bottomrule
\end{tabular}
\end{table}

\begin{table}[htbp]
\centering
\caption{Pituitary and mediobasal hypothalamic gene expression in LepRb$\Delta$Vglut2 mice. Pituitary genes normalised to \textit{Actb}, MBH genes to \textit{B2m}. ns = not significantly different.}
\label{tab:qpcr}
\begin{tabular}{llcl}
\toprule
Gene & Tissue & Result & Statistic \\
\midrule
\textit{Gnrhr} & Pituitary & ns & $t_{10}=0.86$, $p=0.41$ \\
\textit{Lh$\beta$} & Pituitary & ns & $t_{10}=1.47$, $p=0.17$ \\
\textit{Fsh$\beta$} & Pituitary & ns & $t_{9}=1.63$, $p=0.14$ \\
\textit{Cga} & Pituitary & ns ($100\pm17.47\%$ vs $245.3\pm89.91\%$) & $t_{8}=1.20$, $p=0.26$ \\
\textit{Kiss1} & MBH & ns & $t_{8}=1.42$, $p=0.19$ \\
\textit{Pdyn} & MBH & ns & $t_{8}=0.02$, $p=0.98$ \\
\textit{Kiss1r} & MBH & ns & $t_{8}=1.44$, $p=0.19$ \\
\textit{Tac3r} & MBH & ns & $t_{8}=0.99$, $p=0.35$ \\
\textit{Tac2} & MBH & $\sim$50\% decrease & $t_{8}=3.39$, $p=0.009$ \\
\bottomrule
\end{tabular}
\end{table}

\begin{table}[htbp]
\centering
\caption{PMv-targeted LepR re-expression with or without Vglut2 deletion in LepR$^{loxTB}$ mice.}
\label{tab:rescue}
\begin{tabular}{lccc}
\toprule
Group & $n$ & Puberty / pregnancy & PMv pSTAT3 / ovary statistic \\
\midrule
Wild-type & 13 & Fertile & baseline (positive control) \\
LepR$^{loxTB}$ (null) & 4 & Infertile & (negative control) \\
AAV-Cre LepR$^{loxTB}$ & 5 & Rescued & PMv pSTAT3: $F_{3,13}=73.53$, $p<0.0001$ \\
AAV-Cre LepR$^{loxTB}$;Vglut2$^{flox}$ & 6 (on-target) & No VO, no pregnancy & Vglut2 ISH: $t_{8}=4.20$, $p=0.003$ \\
 & & & Corpora lutea: $F_{3,17}=31.16$, $p<0.0001$ \\
\bottomrule
\end{tabular}
\end{table}

% ==========================================================================
% Figures (from figures/figures.tex)
% ==========================================================================
\begin{figure}[htbp]
\centering
\includegraphics[width=0.95\textwidth]{figures/fig1_design.pdf}
\caption{Experimental design for chemogenetic activation of PMv LepRb neurons. Adult LepRb-Cre females (8--10 weeks) received unilateral stereotaxic injection of $\sim$50~nL of AAV-hM3Dq-mCherry into the ventral premammillary nucleus (PMv). Four weeks later, animals were connected to a Culex automated blood-sampling system; three baseline samples were drawn at 10-min intervals, and intravenous clozapine-$N$-oxide (CNO) or clozapine (0.5~mg/kg in 50~$\mu$L) was delivered at time 0. Serial LH measurements continued every 10~min for 60~min; brains were processed for cFos and mCherry immunofluorescence 2~h post-injection.}
\label{fig:design}
\end{figure}

\begin{figure}[htbp]
\centering
\includegraphics[width=0.7\textwidth]{figures/fig2_responders.pdf}
\caption{Proportion of mice with an LH rise greater than 0.5~ng/mL within 10--20~min of intravenous ligand injection, grouped by AAV targeting and ligand: 8/15 (55\%) of CNO-injected PMv-hit animals, 4/4 (100\%) of clozapine-injected PMv-hit animals, 0/7 of PMv-miss plus AAV-mCherry controls (one animal showed a rise at 30~min, after the defined window), and 2/7 (28.5\%) of no-AAV animals given clozapine alone. Absolute responder counts are printed above each bar.}
\label{fig:responders}
\end{figure}

\begin{figure}[htbp]
\centering
\includegraphics[width=0.5\textwidth]{figures/fig3_cycles.pdf}
\caption{Estrous cycles completed per 30 days in adult LepRb-Cre control ($4.47\pm0.30$) and LepRb$\Delta$Vglut2 ($3.31\pm0.36$) females. Bars show mean $\pm$ SEM; unpaired Student's $t$-test, $t_{26}=2.45$, $p=0.021$.}
\label{fig:cycles}
\end{figure}

\begin{figure}[htbp]
\centering
\includegraphics[width=0.75\textwidth]{figures/fig4_mbh_qpcr.pdf}
\caption{Mediobasal hypothalamic qPCR for KNDy-network transcripts in Vglut2$^{flox}$ ($n=5$) and LepRb$\Delta$Vglut2 ($n=5$) females, normalised to $\beta$-2-microglobulin. Bars show $-\log_{10}(p)$ for each comparison; dashed line marks $p=0.05$. Only \textit{Tac2} differed significantly between genotypes ($t_{8}=3.39$, $p=0.009$).}
\label{fig:mbhqpcr}
\end{figure}

\begin{figure}[htbp]
\centering
\includegraphics[width=0.85\textwidth]{figures/fig5_model.pdf}
\caption{Working model. Leptin acts on LepRb neurons of the PMv, most of which are glutamatergic (Vglut2$^{+}$). These neurons provide excitatory drive to GnRH neurons directly and, indirectly, to the Arc KNDy network through glutamate. In LepRb$\Delta$Vglut2 mice, the glutamatergic arm is disabled; downstream, hypothalamic \textit{Tac2} expression falls by $\sim$50\% and the LH response to exogenous kisspeptin-10 is blunted (AUC $t_{11}=3.14$, $p=0.009$).}
\label{fig:model}
\end{figure}

\begin{figure}[htbp]
\centering
\includegraphics[width=0.7\textwidth]{figures/fig6_rescue.pdf}
\caption{Qualitative summary of leptin-induced pSTAT3 immunoreactivity in the PMv across the four groups of the rescue experiment: wild-type, LepR$^{loxTB}$ null, AAV-Cre-injected LepR$^{loxTB}$, and AAV-Cre-injected LepR$^{loxTB}$;Vglut2$^{flox}$. Ordinal heights indicate that PMv pSTAT3 is restored to wild-type levels by AAV-Cre in both injected genotypes, while uninjected LepR$^{loxTB}$ mice show baseline values (overall ANOVA $F_{3,13}=73.53$, $p<0.0001$). Despite restored LepR signalling, LepR$^{loxTB}$;Vglut2$^{flox}$ females with correct PMv injections showed no vaginal opening or pregnancy.}
\label{fig:rescue}
\end{figure}

\end{document}
